\section{Conclusion}
\label{sec:conclusion}

This paper presented a physics-guided two-stage XGBoost workflow for 5-min-grid PV plant power prediction. The method separates atmospheric transmissivity estimation from final AC power prediction, giving the model a physically meaningful intermediate representation while retaining the practical advantages of tabular gradient boosting. On the chronological daylight test split, under an offline target-weather-input protocol, the proposed path achieved an RMSE of \RMSEproposed{}, an nRMSE of \NRMSEproposed{}, and an $R^2$ of \RtwoProposed{}. Ablation results showed that the $\kt$ intermediate improves the forecast relative to direct single-stage XGBoost, while cloud-layer separation gives a smaller but measurable gain over total cloud cover alone.

The value of the approach is not that XGBoost is new, but that the prediction workflow is organized around solar physics, plant operation, and a clear deployment template. The baseline study narrows the claim: the model improves over previous-day same-time persistence, while a live 5-min last-value nowcast remains stronger under the current information setting. The diagnostic results show where the model is dependable and where weather-driven ramps remain difficult. Actual operational forecasting claims should be tested next with archived issue-time weather forecasts at explicit lead times such as 5 min, 30 min, 1 h, and day-ahead.
